\chapter{Integrated Results and Comparative Analysis}
\label{ch:results}

Chapter~\ref{ch:activities} reconstructed the project chronologically, phase
by phase. This chapter does not repeat that chronology. It instead
integrates the evidence across phases to answer the scientific questions
that recur throughout the project's history: how reliably visible corrosion
and hidden structural condition can each be estimated, what images
contribute once metadata is already available, how conclusions change under
stricter evaluation, what can and cannot be said about degradation and
service life, and what the classification branch currently supports. Facts
already established in Chapters~\ref{ch:dataset}--\ref{ch:activities} are
cited by reference rather than re-derived.

\section{Scope and Rules for Cross-Phase Comparison}
\label{sec:results-scope}

The project's successive phases differ in dataset basis (792 rows with one
imputed observation versus 791 aligned rows, Section~\ref{sec:dataset-audit}),
target definition, feature representation, split construction, number of
folds, model-selection procedure, use of manually assigned labels, and, for
the categorical targets, the number of severity classes
(Table~\ref{tab:category-schema}). A lower MAE or higher $R^2$ in one phase
than another does not, by itself, establish that one phase's approach is
superior: it may instead reflect a different, non-comparable protocol.

This chapter therefore distinguishes three kinds of comparison: a
\emph{direct numerical comparison}, permitted only where target, feature
set, dataset basis, and evaluation regime are sufficiently aligned, with
remaining protocol differences stated alongside the numbers; a \emph{ratio
comparison}, which normalises a result against its own phase's baseline
(e.g.\ a stricter-regime MAE divided by that same model's grouped-holdout
MAE) so only relative degradation, not absolute error, crosses phases; and
a \emph{qualitative comparison} of scientific conclusions, robustness
patterns, and failure modes rather than metric values. Most comparisons
below are of the second or third kind.

\section{Visible-Corrosion Estimation}
\label{sec:results-visible}

\subsection{Current-corrosion estimation}

Across the project, visible surface corrosion is the quantity estimated
with the most consistent success. The classical baseline reached a grouped-
holdout $R^2$ of 0.900 for current-week peak rust
(Section~\ref{sec:phase1}); the interpretable multi-family feature phase
reached $R^2=0.646$ (total rust) and $R^2=0.784$ (peak rust) under the same
regime (Section~\ref{sec:phase3-corrosion}). These results are not placed
on a single ranked list, because they use different feature sets and, in
the classical baseline's case, a dataset basis and label-construction
method that materially affects how the result should be read (next
subsection). Taken together, however, they support a first, qualified
conclusion: image-derived representations of visible corrosion carry
genuine, learnable signal under grouped-specimen evaluation.

\subsection{Interpretability and the tautology risk}

The strength of this signal must be read against two distinct label-
reconstruction concerns, which implicate different targets and different
degrees of severity and should not be merged into one claim. In the
classical baseline, the current-corrosion (peak-rust) prediction is
dominated (91.5\% of feature importance) by a single rust-mask-derived
colour-threshold feature closely aligned with the target
(Section~\ref{sec:phase1}); this is a substantial shortcut/tautology risk,
but the available evidence establishes only that this project-internal
feature is closely aligned with the manually assigned label, not that it
is computed by literally the same rule used to assign that label. In the
robustness-quantification phase, by contrast, the near-identity is
directly quantified: \texttt{surface\_total\_rust\_pct} is correlated with
a separate image feature at Spearman $\rho=0.9996$
(Section~\ref{sec:phase4a}) — close to near label reconstruction, a
stronger claim than a shortcut risk. These two findings concern different
targets (peak rust in the classical baseline; total/surface rust in the
robustness-quantification phase), so they illustrate a recurring class of
risk — label-adjacent features dominating a benchmark — rather than two
data points about the same target. Peak rust in the later phases, for
which no comparably near-tautological pairing has been found, remains the
more meaningful test of genuine image-driven estimation, and this chapter
treats these label-adjacent benchmarks as supportive rather than as
headline evidence on their own.

\subsection{Robustness}

Robustness to distribution shift is where the picture becomes more
qualified still, and where MAE, $R^2$, and Spearman do not always tell the
same story — and where the two phases that tested this do not agree with
each other for treatment-level shift. In the interpretable-feature phase,
the best-available regime model for each target remains positive under
leave-one-campaign-out ($R^2=0.582$ and $0.595$) but turns negative under
leave-one-treatment-out ($R^2=-0.677$ and $-0.077$;
Table~\ref{tab:phase3-results}), evaluated over seven coarse treatment
protocols (Section~\ref{sec:phase3-structural}) — a genuine collapse in
this phase's own treatment-holdout evaluation.

The robustness-quantification phase's own leave-one-treatment-out
evaluation does \emph{not} reproduce this fragility. Tracking each
target's own grouped-holdout-winning model under every regime — the
appropriate same-model comparison, and how this phase's own robustness
table is itself constructed — its leave-one-treatment-out MAE is at or
\emph{below} its grouped-holdout MAE for both targets (ratios 0.55 and
0.76), mean $R^2$ is 0.995 (total rust) and 0.957 (peak rust), and Spearman
correlation stays above 0.97 throughout, computed directly this session
from the phase's own saved per-fold outputs
(Section~\ref{sec:phase4a}). This phase's leave-one-treatment-out
evaluation also uses nine, more finely grouped treatment splits rather
than the earlier phase's seven coarse protocols, which may partly explain
the difference; the evidence available here establishes the outcome, not
the mechanism.

Leave-one-campaign-out robustness, by contrast, does replicate: both
phases keep $R^2$ positive (or an MAE ratio at or below 1.0) for both
visible-corrosion targets. Visible-corrosion estimation should therefore
be described as robust to campaign-level shift in both phases; treatment-
level robustness is a genuine failure mode in one phase but is not
established as a project-wide limitation, since the other phase's own
treatment-holdout evaluation shows no comparable deterioration on any of
the three metrics.

\section{Structural-Condition Prediction}
\label{sec:results-structural}

Under a grouped specimen holdout — the least demanding of the three
regimes — hidden structural condition shows a modest but non-trivial
relationship to the available features: wire-area loss reaches
$R^2=0.310$ and ultimate load reaches $R^2=0.385$ in the interpretable-
feature phase (Section~\ref{sec:phase3-structural}). This grouped-holdout
result looks more encouraging than the strict-transfer results below for a
specific, verified reason: it evaluates a random subset of specimens drawn
from the same experimental mixture the model was trained on
(Section~\ref{sec:grouped-holdout}), whereas the stricter regimes test
extrapolation across specimen design or treatment protocol.

Under leave-one-campaign-out, structural performance collapses in both
phases that tested it: mean $R^2$ falls to $-0.39$ (wire-area loss) and
$-8.36$ (ultimate load) in the interpretable-feature phase, and the
independently re-audited robustness benchmark shows ultimate load's
leave-one-campaign-out MAE inflating to 3.24 times its own grouped-holdout
value, with Spearman falling to 0.35; wire-area loss's leave-one-campaign-out
Spearman turns negative ($-0.089$), worse than no rank information at all
(Section~\ref{sec:phase4a}). Within a single mesh family, the same
collapse appears a third way, without invoking either stricter regime at
all: every feature-set configuration tested for terminal ultimate load
achieves a negative mean $R^2$ inside the 4-mesh or 7-mesh group considered
alone, even where the raw MAE looks numerically reasonable
(Section~\ref{sec:phase4b}). Spearman correlation is what exposes this
gap in every case: a model or a pooled correlation can report a
plausible-looking absolute error while its rank correlation reveals that it
carries little or no genuine discriminative signal once the sample is
restricted to a design-homogeneous subgroup.

Two feature-importance diagnostics, computed independently in two phases,
converge on the same explanation. The interpretable-feature phase's
full-data ultimate-load model is dominated by steel-mesh count and
campaign indicators, not image descriptors
(Section~\ref{sec:phase3-structural}); the robustness-quantification
phase's five strongest univariate correlates of ultimate load are five
metadata/temporal variables tied at exactly $|r|=0.827$, a direct signature
of their perfect alignment with campaign identity on the 48-row structural
subset (Section~\ref{sec:confounding}, Section~\ref{sec:phase4a}). This
sparse-supervision constraint (48 rows, one per specimen;
Section~\ref{sec:structural-supervision}) is the most plausible reason
structural inference is systematically weaker than visible-corrosion
inference — a difference in available supervision, not primarily in
modelling effort.

A further distinction matters for how the project's own best evidence on
this question should be read. The interpretable-feature phase's structural
models are mixed-input — metadata, manually assigned corrosion labels, and
image descriptors together (Section~\ref{sec:phase3-structural}) — whereas
the robustness-quantification phase's \emph{final}, robustness-weighted
structural model selection uses metadata alone for both targets: a
RandomForest on metadata-only features for ultimate load, and a CatBoost on
metadata-only features for wire-area loss (Section~\ref{sec:phase4a}). The
project's own best-performing, most carefully selected structural models
therefore use no image input at all. This is stronger and more specific
than saying structural inference is weak: it is direct evidence about where
whatever usable structural signal remains actually originates.

\emph{Bounded conclusion:} the available evidence supports substantially
weaker and less stable structural inference than visible-corrosion
inference, driven by sparse supervision and campaign-linked confounding,
and the project's own final, robustness-weighted structural models do not
use image input. This does not establish that structural prediction from
corrosion images is impossible in principle — only that it is not
established by the evidence gathered here.

\section{Incremental Value of Image Information}
\label{sec:results-incremental}

A distinction the preceding two sections make necessary, and which this
section addresses directly, is between two claims that are easily
conflated: that images contain predictive signal, and that images add
incremental predictive value once metadata is already available. The first
is well supported for visible-corrosion estimation (Section~\ref{sec:results-visible}).
The second requires a design that isolates the image contribution from the
metadata contribution — an image-only, metadata-only, and combined
comparison under matched conditions (Section~\ref{sec:multimodal}) — and
this design was not applied consistently across the project's history.

The deep-embedding phase compared only an image-only and a combined
image-plus-metadata configuration; no metadata-only baseline was trained
(Section~\ref{sec:phase2}, Section~\ref{sec:multimodal}). Its result — a
lower validation error but a higher test error for the combined model
relative to the image-only model — can therefore speak only to whether
adding metadata to an image model helped in that one saved run; it cannot
answer whether images add value beyond metadata, because the metadata-only
anchor needed to answer that question was never computed. This chapter does
not treat that phase as if it had run the complete ablation.

The terminal-load-focused phase is the only phase that made the full
three-way comparison a designed, central part of its methodology
(Section~\ref{sec:phase4b}). Its pooled result, across all 48 specimens, is
unambiguous: a metadata-only Ridge regression achieves the lowest MAE
(0.173~kN) and the second-highest Spearman correlation (0.758) of every
configuration tested, and adding RGB descriptors measurably worsens both
MAE and $R^2$ (Table~\ref{tab:phase4b-featuresets}). A subset restricted to
specimens with measurable terminal corrosion shows a narrower, real
exception — a combined feature set edges ahead of metadata-only on MAE and
$R^2$ there — but metadata-only still holds a marginally higher Spearman
correlation even within that subset, and the phase's own primary analysis
is the pooled one, not the subset one (Section~\ref{sec:phase4b}). This
subset result is genuine evidence that images become more informative once
corrosion has visibly progressed; it is not evidence that overturns the
pooled conclusion.

\emph{Bounded conclusion:} for terminal ultimate load — the only target
with a complete three-way ablation — metadata is the strongest available
pooled predictor, and visible corrosion descriptors do not provide a
reliable pooled improvement over it, with a narrow, non-overturning
exception once corrosion has visibly onset. That images carry some signal
is not the same statement as that images add incremental value beyond
metadata for this target, and the evidence supports the first claim more
than the second.

\section{Robustness and Distribution Shift}
\label{sec:results-robustness}

Section~\ref{sec:eval-regimes} defines the three evaluation regimes; this
section asks what changed, across the project, when evaluation became
harder — and where the two phases that tested more than one regime agree
or disagree with each other. Leave-one-campaign-out is the regime that
replicates most cleanly: structural targets are fragile under it in both
phases, and visible-corrosion targets remain robust under it in both
phases (Sections~\ref{sec:results-visible},~\ref{sec:results-structural}).
Leave-one-treatment-out is where the two phases diverge for
visible-corrosion targets specifically: the interpretable-feature phase's
visible-corrosion models collapse under it, while the robustness-
quantification phase's do not
(Section~\ref{sec:results-visible}) — a phase-specific finding, not a
replicated project-wide pattern. Structural targets, by contrast, are
fragile under both stricter regimes in both phases, with
leave-one-campaign-out typically the more damaging of the two — consistent
with campaign identity, not treatment identity, carrying most of the
design-linked confounding that structural models draw on
(Section~\ref{sec:confounding}, Section~\ref{sec:results-structural}).
Because campaign holdout tests extrapolation across a design space in
which mesh configuration, chloride level, and ageing duration all shift
together (Section~\ref{sec:confounding}), while treatment holdout tests
transfer within one design point to an unseen protocol, campaign-level
fragility being the structural layer's dominant weak point is a plausible
consequence of \emph{what} that regime perturbs.

\begin{figure}[htbp]
\centering
\includegraphics[width=0.62\textwidth]{fig_5_1_robustness_synthesis.png}
\caption{Relative MAE (regime MAE divided by that same model's own
grouped-holdout MAE) for the two visible-corrosion and two structural
targets, using each phase's own grouped-holdout-winning model evaluated
under every regime — a fixed-model comparison, not a
best-available-model-per-regime one.\footnote{Fixed models: the
interpretable-feature phase uses RandomForest for total rust, wire-area
loss, and ultimate load, and an MLP for peak rust; the
robustness-quantification phase uses GradientBoosting for surface total
rust and RandomForest for the remaining three targets. The
robustness-quantification phase's benchmark table is computed on this
fixed-model basis by its own diagnostics code; the interpretable-feature
phase's values were recomputed this session using the same rule against
its own saved per-fold metrics.} No raw MAE from one phase is compared
against another phase's raw MAE. Structural targets inflate under both
regimes in both phases, most severely for ultimate load under
leave-one-campaign-out; visible-corrosion targets stay close to or below
their own grouped-holdout MAE under leave-one-campaign-out in both phases,
but diverge under leave-one-treatment-out — fragile in the
interpretable-feature phase, robust in the robustness-quantification
phase.}
\label{fig:results-robustness-synthesis}
\end{figure}

Figure~\ref{fig:results-robustness-synthesis} shows both the replication
and the divergence directly: the structural-target inflation pattern and
the visible-corrosion targets' campaign-level robustness recur
independently across two phases with unrelated feature engineering, which
makes it unlikely that these two patterns are artefacts of a single
phase's feature-engineering choices; the visible-corrosion treatment-level
fragility does not recur, and should be read as specific to the
interpretable-feature phase's own treatment-holdout evaluation rather than
as a property of the data itself.

\section{Degradation and Proxy-RUL Evidence}
\label{sec:results-degradation}

Two generations of the degradation/proxy-RUL derivation
(Section~\ref{sec:proxy-rul}) exist in the project, on configurations that
are not directly comparable. The first
(Section~\ref{sec:phase3-degradation}) fits a composite four-series
derivation — two observed visible-corrosion series and two model-estimated
structural series — against three threshold types over a 104-week
projection horizon: of 48 specimens, 33 receive a finite estimate, 23 of
which are exactly zero weeks, leaving only ten with a proxy-RUL that is
both finite and genuinely forward-looking.

The second generation instead derives threshold-crossing status from a
single series (model-estimated wire-area-loss fraction) against three
absolute thresholds over a shorter, 365-day horizon — a narrower target
and a materially different horizon, so the two are compared here only at
the level of their conclusions, not their counts. Its own saved output
classifies every one of the 144 specimen/threshold combinations tested as
already past the threshold at the first observation, crossed during
observation, or not crossed within the horizon: there is no genuinely
forward-looking finite crossing anywhere in this refined output, a still
more conservative result than the first generation's ten forward-looking
cases. The project's own documentation separately flags, as an unresolved
validation concern rather than a confirmed defect, that most specimens
already show a model-estimated wire-area loss at or above the lowest
threshold at their first observation — consistent with most crossings
being attributed to the baseline rather than to projected degradation.

\emph{Bounded conclusion:} both generations of the degradation/proxy-RUL
pipeline demonstrate that the derivation chain defined in
Section~\ref{sec:proxy-rul} can be executed end to end, and both
generations' own outputs argue against treating the resulting numbers as
validated remaining-life estimates: the first because most of its finite
estimates are degenerate, the second because it produces no forward-looking
estimate at all under a more conservative configuration. True remaining
useful life remains unavailable in this dataset regardless of which
generation is consulted.

\section{Current Four-Class Classification Branch}
\label{sec:results-classification}

The four-class classification effort (Section~\ref{sec:phase5}) is
structurally independent of every modelling question addressed above: it
targets a different task (discrete severity classification rather than
continuous estimation), a different, current four-level schema
(Table~\ref{tab:category-schema}), and shares no code with any earlier
phase. Data preparation for this task is complete and verified: the target
schema, the class-imbalance characterisation, the specimen-level grouped
split, the leakage-safety runtime check, and the augmentation pipeline all
exist and have been run. No classifier has been trained, compared, or
evaluated on this data at the time of writing. This branch therefore
contributes methodology and data-preparation evidence to this report, not
model-performance evidence, and it is not compared against the
interpretable-feature phase's five-category classification results
(Section~\ref{sec:phase3-corrosion}), which model a different, historical
schema with a different number of classes.

\section{Integrated Findings}
\label{sec:results-integrated}

Table~\ref{tab:cross-phase-comparability} summarises which quantitative
results in this project may be compared directly, which may only be
compared as ratios to their own baseline, and which must be compared only
qualitatively. Table~\ref{tab:integrated-findings} then states, for each
scientific question addressed in this chapter, the project's current
evidence level, main finding, main limitation, and the most defensible
claim the evidence supports. Evidence-level terms are used as defined in
the surrounding text of this chapter and are not numeric scores: \emph{well
supported} indicates a conclusion replicated across phases and stable under
at least one stricter evaluation regime; \emph{conditionally supported}
indicates a conclusion that holds under some but not all tested conditions;
\emph{exploratory} indicates a working demonstration not yet validated
against an independent standard; \emph{not yet evaluated} indicates that no
model-performance evidence exists at all.

\begin{table}[htbp]
\centering
\caption{Cross-phase comparability. What may be compared directly, as a
ratio to each result's own baseline, or only qualitatively, and why.}
\label{tab:cross-phase-comparability}
\small
\begin{tabular}{p{3.2cm}p{3.3cm}p{6.0cm}}
\toprule
Comparison & Permitted form & Reason \\
\midrule
Interpretable-feature phase vs.\ robustness-quantification phase, same
target family, robustness pattern & Ratio (regime MAE / own
grouped-holdout MAE) or qualitative & Different feature sets and dataset
bases (792 vs.\ 791 rows); absolute MAE not comparable, but both phases
independently test the same robustness question. \\
Classical baseline vs.\ interpretable-feature phase, current-corrosion
accuracy & Qualitative only & Different feature engineering, different
label-construction tautology risk (Section~\ref{sec:results-visible}); not
placed on a shared ranked list. \\
Interpretable-feature phase vs.\ terminal-load-focused phase, ultimate-load
MAE & Not compared numerically & Different feature sets, splits, and
dataset bases (Section~\ref{sec:phase3-structural}); comparing 0.121~kN
against 0.174~kN as a leaderboard would be a category error. \\
Deep-embedding phase vs.\ terminal-load-focused phase, image value & Not
compared numerically & Different targets (current-week peak rust vs.\
terminal ultimate load) and different representations; deep-embedding
phase lacks a metadata-only baseline (Section~\ref{sec:results-incremental}). \\
First-generation vs.\ second-generation proxy-RUL & Qualitative (both
support the same bounded conclusion) & Different target composition,
threshold definitions, and projection horizons (104 weeks $\approx$ 728
days versus 365 days $\approx$ 52 weeks;
Section~\ref{sec:results-degradation}). \\
Historical five-class classification (interpretable-feature phase) vs.\
current four-class classification & Not compared & Different category
schema (Table~\ref{tab:category-schema}); the current branch has no
trained-model result to compare. \\
\bottomrule
\end{tabular}
\end{table}

\FloatBarrier

\begin{table}[htbp]
\centering
\caption{Integrated scientific findings across the project. Evidence-level
terms are defined in the surrounding text and are not numeric scores.}
\label{tab:integrated-findings}
\footnotesize
\begin{tabular}{p{2.5cm}p{1.9cm}p{3.3cm}p{3.3cm}}
\toprule
Question & Evidence level & Main finding & Defensible claim \\
\midrule
Visible surface corrosion (Section~\ref{sec:results-visible}) & Well
supported (peak rust; campaign-level robustness); conditionally supported
(total rust, tautology risk; treatment-level robustness, phase-dependent)
& Robust under campaign holdout in both phases; treatment-holdout
robustness diverges between phases (fragile in one, robust in the other)
& Images carry genuine, learnable visible-corrosion signal, most credibly
for peak rust, robust under campaign-level transfer; treatment-level
transfer is unresolved, not established as a project-wide weakness \\
Hidden structural condition (Section~\ref{sec:results-structural}) &
Conditionally supported, weak & Grouped-holdout signal collapses under
both stricter regimes; final selected models are metadata-only & The
evidence supports substantially weaker structural inference than visible-
corrosion inference; not established that images cannot help in
principle \\
Terminal ultimate load / incremental image value
(Section~\ref{sec:results-incremental}) & Conditionally supported (against
image value) & Metadata-only is the strongest pooled predictor; a narrow
post-onset exception exists & Images do not add a reliable pooled
improvement over metadata for terminal ultimate load \\
Cross-regime / cross-campaign robustness (Section~\ref{sec:results-robustness})
& Well supported (campaign-level pattern; structural fragility);
conditionally supported (treatment-level visible-corrosion behaviour,
phase-dependent) & Structural fragility under both stricter regimes, and
visible-corrosion robustness under campaign holdout, replicate across two
phases; treatment-holdout behaviour for visible-corrosion diverges between
phases & Robustness behaviour is target-family- and regime-specific;
campaign-level and structural findings replicate across phases, but
treatment-level visible-corrosion robustness does not and should be
treated as phase-dependent \\
Degradation / proxy-RUL (Section~\ref{sec:results-degradation}) &
Exploratory & Derivation chain executable end to end in two non-comparable
configurations; both argue against forward-looking confidence & Useful as
an exploratory screening framework; not a validated remaining-life
estimate \\
Four-class classification (Section~\ref{sec:results-classification}) & Not
yet evaluated & Data preparation, splitting, and augmentation complete; no
classifier trained & Contributes methodology evidence only; no
performance claim can yet be made \\
\bottomrule
\end{tabular}
\end{table}

\FloatBarrier

Read together, these rows describe a project whose strongest evidence
concerns visible-corrosion estimation, reinforced across independent
phases except for treatment-level robustness, which is phase-specific
rather than replicated (Section~\ref{sec:results-robustness}). Its most
consequential negative findings were established in different ways and
should not be flattened into one narrative of replication: campaign-
confounded structural supervision was named early and reinforced by later,
independent requantification (Sections~\ref{sec:results-structural}
and~\ref{sec:phase4a}); proxy-RUL's non-validated status is reinforced by
two independently configured generations of the same derivation
(Section~\ref{sec:results-degradation}); but the absence of reliable
incremental image value for terminal ultimate load rests on one phase's
own designed three-way ablation
(Section~\ref{sec:results-incremental}), not on independent replication by
a separate experiment, and is best read as a single well-designed test
rather than a cross-phase-confirmed pattern.
